\section{Introdução}

O planejamento do transporte de cargas constitui um tema clássico da Pesquisa
Operacional, com desdobramentos em problemas de consolidação, designação de
veículos, dimensionamento de frota e planejamento operacional de embarques.
Revisões amplas, como as de \textcite{crainic1997planning} e
\textcite{steadieseifi2014multimodal}, mostram a extensão desse campo e deixam
claro que diferentes configurações logísticas geram estruturas decisórias
próprias. Nesse contexto, o presente trabalho parte de uma revisão rápida da
literatura para posicionar um problema logístico específico que, embora
relacionado a decisões de alocação de frota, apresenta uma combinação de
restrições operacionais ainda pouco explorada como objeto central de modelagem.

A justificativa para essa revisão rápida é direta. Antes de aprofundar a
formulação matemática do problema, é necessário verificar como a literatura vem
tratando decisões de atribuição entre cargas e veículos em contextos nos quais
há restrições de sequenciamento físico, heterogeneidade de frota e impacto
econômico associado ao uso parcial da capacidade dos equipamentos. Assim, a
revisão não constitui o produto final do estudo, mas sim uma etapa de
delimitação conceitual: ela orienta a definição do problema, ajuda a identificar
lacunas e sustenta a transição para a etapa principal da pesquisa, centrada no
desenvolvimento e na análise de um modelo de Pesquisa Operacional.

A motivação empírica inicial surgiu em operações de distribuição terrestre de
bobinas de aço importadas no Brasil. De acordo com o \textcite{acobrasil2025},
o volume de importações siderúrgicas cresceu de forma expressiva nos últimos
anos, e uma parcela relevante desses fluxos é posteriormente distribuída por
transporte rodoviário. Em muitas operações, as bobinas são descarregadas em
portos e removidas para recintos alfandegados secundários, onde permanecem
armazenadas até a nacionalização e a expedição. Nesses ambientes, o arranjo
físico da carga no pátio ou armazém, frequentemente associado à ordem de
recebimento ou ao processo de armazenagem, induz uma sequência prática de acesso
às unidades. A retirada fora dessa ordem tende a demandar rehandling, elevando
tempo, custo e risco operacional.

Essa configuração leva a uma decisão operacional relevante: dada uma sequência
de unidades disponíveis para expedição e uma frota heterogênea, como atribuir
subsequências contíguas de cargas aos veículos disponíveis? Em muitos contratos
de frete, essa decisão possui consequências econômicas diretas, pois os valores
pagos dependem não apenas da tonelagem embarcada, mas também de mínimos
faturáveis por tipo de veículo. Assim, duas atribuições distintas podem mover o
mesmo volume total e, ainda assim, produzir custos operacionais bastante
diferentes.

Ao longo do amadurecimento do problema, no entanto, tornou-se claro que a
motivação por bobinas de aço descreve apenas um caso particular de uma classe
mais geral. O elemento estrutural decisivo não é o tipo de mercadoria em si,
mas a existência de cargas para as quais, uma vez atribuída uma subsequência a
um determinado veículo, existe uma única forma operacionalmente válida de
carregamento. Em outras palavras, o problema relevante não inclui um subproblema
combinatório adicional de arranjo interno da carga no veículo: esse arranjo é
univocamente determinado pelas características da carga, pela configuração do
equipamento e pelas regras operacionais do contexto.

Essa observação amplia o alcance da pesquisa. O caso das bobinas continua sendo
um motivador importante, mas o modelo pretendido busca representar uma classe de
operações logísticas em que: (i) as unidades de carga estão disponíveis em uma
sequência física de acesso; (ii) a expedição é feita por uma frota heterogênea;
(iii) a decisão central consiste em atribuir subsequências contíguas de carga a
veículos; e (iv) o carregamento em cada veículo é unicamente determinado depois
da atribuição. O foco, portanto, desloca-se de uma revisão sistemática sobre um
tema setorial para uma pesquisa orientada a modelo, informada por uma revisão
rápida da literatura e fundamentada em uma estrutura decisória de interesse mais
geral para a Pesquisa Operacional aplicada à logística.

Com essa delimitação, este trabalho passa a perseguir dois objetivos
complementares. O primeiro é organizar, por meio de uma revisão rápida, o que a
literatura oferece de mais diretamente relacionado a problemas de alocação de
frota, atribuição de veículos e planejamento de transporte de cargas com
restrições operacionais próximas ao contexto estudado. O segundo, que se torna o
objetivo principal da pesquisa, é formalizar o problema logístico descrito,
evidenciar suas particularidades e preparar o terreno para a construção e a
análise de formulações matemáticas coerentes com essa estrutura decisória. Como
desdobramento desse segundo objetivo, o trabalho também inclui a implementação
de um protótipo computacional da formulação completa e a preparação de uma base
de dados de execuções, destinada a sustentar a etapa empírica posterior da
pesquisa. A seção seguinte apresenta essa transição de modo explícito, primeiro
por meio de uma formulação inteira do tipo \emph{set partitioning over
intervals} e, em seguida, por meio de uma versão simplificada em programação
dinâmica que ajuda a explicitar a lógica sequencial do problema e seus limites.
